\section*{Chapitre \ref{ch:b3:point-groups} --- Symétrie et groupes ponctuels}

\begin{solution}{exo:b3:point-groups:1}
Éthène~: trois $C_2$ perpendiculaires, trois plans, $i$~: $D_{2h}$. \ce{XeF4} (plan
carré)~: $C_4$, quatre $C_2 \perp C_4$, $\sigma_h$, $2\sigma_v$, $2\sigma_d$, $i$, $S_4$~:
$D_{4h}$. \ce{PCl5} (bipyramide trigonale)~: $C_3$, trois $C_2$, $\sigma_h$, $3\sigma_v$,
$S_3$~: $D_{3h}$. 1,3,5-Trichlorobenzène~: les mêmes éléments que \ce{BF3}~: $D_{3h}$.
\end{solution}

\begin{solution}{exo:b3:point-groups:2}
Éthane décalé~: $C_3$, trois $C_2$ perpendiculaires, trois $\sigma_d$, $i$, $S_6$~:
$D_{3d}$. Éclipsé~: $C_3$, trois $C_2$, $\sigma_h$, trois $\sigma_v$~: $D_{3h}$.
Cyclohexane chaise~: $D_{3d}$.
\end{solution}

\begin{solution}{exo:b3:point-groups:3}
\ce{SO3} ($D_{3h}$)~: non. \ce{SO2} ($C_{2v}$)~: oui. \ce{PCl3} ($C_{3v}$)~: oui.
\ce{XeF4} ($D_{4h}$)~: non. \ce{CH2Cl2} ($C_{2v}$)~: oui. cis-1,2-dichloroéthène
($C_{2v}$)~: oui~; trans ($C_{2h}$)~: non.
\end{solution}

\begin{solution}{exo:b3:point-groups:4}
$C_2(z) = \mathrm{diag}(-1,-1,1)$, $\chi = -1$~; $i = \mathrm{diag}(-1,-1,-1)$,
$\chi = -3$~; $\sigma_h(xy) = \mathrm{diag}(1,1,-1)$, $\chi = 1$.
\end{solution}

\begin{solution}{exo:b3:point-groups:5}
Avec les matrices diagonales de l'exercice précédent~: $C_2i = \sigma_h$, $C_2\sigma_h =
i$, $i\sigma_h = C_2$, chaque opération est son propre inverse, et tous
les produits commutent. Le groupe est commutatif, donc $X^{-1}AX = A$ pour
tout $X$~: chaque opération forme une classe à elle seule (quatre
classes, quatre \omterm{def:b3:point-groups:irrep}{représentations irréductibles} de dimension 1).
\end{solution}

\begin{solution}{exo:b3:point-groups:6}
Prenons pour $\sigma_v(1)$ le plan $xz$. Le point $(1,0)$ va par $C_3$ en $(-\frac12,
\frac{\sqrt3}2)$, puis par $\sigma_v(1)$ en $(-\frac12,-\frac{\sqrt3}2)$~;
dans l'autre ordre, il va en $(1,0)$ puis en $(-\frac12,\frac{\sqrt3}2)$.
Les deux produits sont des réflexions par rapport à deux plans différents
($\sigma_v(3)$ et $\sigma_v(2)$)~: le groupe n'est pas commutatif, et
c'est pourquoi ses réflexions forment une classe de trois.
\end{solution}

\begin{solution}{exo:b3:point-groups:7}
Le biphényle tordu conserve le $C_2$ porté par la liaison entre les
cycles et deux $C_2$ qui lui sont perpendiculaires, bissecteurs des
angles entre les plans des cycles~; tous les plans et le \omterm{def:b3:point-groups:inversion}{centre
d'inversion} des formes plane ou perpendiculaire sont perdus. Avec trois
$C_2$ perpendiculaires et aucun élément impropre, le groupe est
$D_2$~: la molécule est chirale (ses deux formes tordues sont des
énantiomères, séparables quand des substituants ortho encombrants
empêchent la rotation).
\end{solution}

\begin{solution}{exo:b3:point-groups:8}
$d_{z^2}$ et $d_{x^2-y^2}$~: $a_1$ ($x^2$, $y^2$, $z^2$ sont $A_1$)~; $d_{xy}$~: $a_2$~; $d_{xz}$~:
$b_1$~; $d_{yz}$~: $b_2$.
\end{solution}

\begin{solution}{exo:b3:point-groups:9}
$1 + 1 + 4 + 9 + 9 = 24 = h$. $\sum g\,\chi_{A_1}\chi_{T_2} = 1\cdot3 + 8\cdot0 + 3(-1)
+ 6(-1) + 6(1) = 0$.
\end{solution}

\begin{solution}{exo:b3:point-groups:10}
Dans le plan perpendiculaire à la droite d'intersection, une réflexion
par rapport à une droite d'angle $\alpha$ envoie l'angle polaire $\phi$
sur $2\alpha - \phi$. Deux réflexions, d'angle
$\alpha$ puis $\alpha + \theta$~: $\phi \to 2\alpha - \phi \to 2(\alpha + \theta) - (2\alpha - \phi) =
\phi + 2\theta$, une rotation de $2\theta$ autour de la droite. Deux plans
verticaux à
$60^\circ$ engendrent donc une rotation de $120^\circ$~: un $C_3$.
\end{solution}

\begin{solution}{exo:b3:point-groups:11}
L'axe C=C=C est un $C_2$ et un $S_4$ (rotation de $90^\circ$, qui échange
les deux plans \ce{CH2}, puis réflexion dans le plan passant par le
carbone central)~; les deux
plans \ce{CH2} sont des $\sigma_d$~; deux $C_2$ perpendiculaires à l'axe
sont leurs bissecteurs.
Ordre 8~: $D_{2d}$, achiral. Dans \ce{ClHC=C=CHCl}, les plans, le $S_4$ et le $C_2$
axial sont détruits par la substitution~; seul subsiste un $C_2$
perpendiculaire à l'axe~: $C_2$, chiral --- un allène à chiralité axiale.
\end{solution}

\begin{solution}{exo:b3:point-groups:12}
\ce{SF6}~: $O_h$ ($h = 48$). \ce{SF5Cl}~: le $C_4$ passant par Cl et quatre $\sigma_v$~:
$C_{4v}$ ($h = 8$). trans-\ce{SF4Cl2}~: $D_{4h}$ ($h = 16$). cis-\ce{SF4Cl2}~: $C_{2v}$ ($h
= 4$). Chacun conserve une partie des opérations de $O_h$, stable par
produit, et 8, 16, 4 divisent 48. Polaires~: \ce{SF5Cl} et
cis-\ce{SF4Cl2}.
\end{solution}

\begin{solution}{pb:b3:point-groups:1}
\textbf{1.} Selon les quatre grandes diagonales passant par C et par
chaque H~; chacune donne
$C_3$ et $C_3^2$~: 8 rotations.
\textbf{2.} Selon $x$, $y$, $z$, par les centres des faces~: $C_2(z)$ envoie $(1,1,1)
\to (-1,-1,1)$, un hydrogène. Une rotation de $90^\circ$ autour de $z$
suivie d'une réflexion dans $xy$ envoie $(1,1,1) \to (-1,1,1) \to (-1,1,-1)$,
un hydrogène~: un
$S_4$~; avec $S_4^3$, 6 opérations.
\textbf{3.} Les plans $x = \pm y$, $y = \pm z$, $x = \pm z$~; $x = y$ contient
$(1,1,1)$ et $(-1,-1,1)$~: chaque plan contient deux liaisons C--H.
\textbf{4.} $1 + 8 + 3 + 6 + 6 = 24$.
\textbf{5.} L'inversion envoie $(1,1,1)$ en $(-1,-1,-1)$, qui n'est pas un
hydrogène~: pas de $i$.
\textbf{6.} $E$~; $8C_3$~; $3C_2$~; $6S_4$~; $6\sigma_d$.
\textbf{7.} \ce{CH3Cl}~: $C_{3v}$, $h = 6$.
\textbf{8.} \ce{CH2Cl2}~: $C_{2v}$, $h = 4$.
\textbf{9.} \ce{CHCl3}~: $C_{3v}$~; \ce{CCl4}~: $T_d$, $h = 24$.
\textbf{10.} \ce{CH2FCl}~: $C_s$ ($h = 2$)~; \ce{CHFClBr}~: $C_1$ ($h = 1$).
\textbf{11.} Chacun conserve les opérations de $T_d$ qui respectent la
substitution~; 6, 4, 6, 24, 2, 1 divisent tous 24.
\textbf{12.} Les trois hydrogènes, permutés par $C_3$, $C_3^2$ et les
trois
$\sigma_v$.
\textbf{13.} Toutes sauf \ce{CCl4}~: selon l'axe $C_3$ pour \ce{CH3Cl} et
\ce{CHCl3}, selon l'axe $C_2$ pour \ce{CH2Cl2}, dans le \omterm{def:b3:point-groups:mirror}{plan de symétrie}
pour \ce{CH2FCl}, dans aucune direction imposée pour \ce{CHFClBr}.
\textbf{14.} \ce{CHFClBr} seule ($C_1$, aucun $S_n$).
\textbf{15.} Deux énantiomères.
\textbf{16.} Ses deux hydrogènes sont équivalents~: le plan passant par
C, F et Cl, bissecteur de H--C--H, est un \omterm{def:b3:point-groups:mirror}{plan de symétrie}.
\textbf{17.} Non~: quatre substituants différents ne laissent que $E$~;
sans $S_n$, la molécule est chirale (et peut être polaire).
\textbf{18.} $(x, y, z)$~: $T_2$.
\textbf{19.} Deux hydrogènes sont les extrémités d'une arête du
tétraèdre~; les 24 opérations envoient toute arête sur n'importe quelle
autre (les 6 arêtes forment un seul ensemble)~: un seul
\ce{CH2Cl2}.
\textbf{20.} \ce{CH2FCl}~: un. \ce{CHFClBr}~: deux, un couple
d'énantiomères.
\textbf{21.} Trois~: ortho ($C_{2v}$), méta ($C_{2v}$), para ($D_{2h}$).
\textbf{22.} Trois~: 1,2,3- ($C_{2v}$), 1,2,4- ($C_s$), 1,3,5- ($D_{3h}$).
\textbf{23.} Deux (atomes de Cl adjacents ou opposés). On ne connaît
qu'un seul \ce{CH2Cl2}, ce qui a exclu un carbone plan et appuyé le
carbone tétraédrique proposé en 1874.
\textbf{24.} $1^2 + 1^2 + 2^2 + 3^2 + 3^2 = 24$~: \textbf{l'ordre de $T_d$,
\omterm{def:b3:point-groups:point-group}{groupe ponctuel} du méthane, est $h = 24$}.
\end{solution}
